\section{Phase II Step 6: Final PoC Metric Selection}
\label{sec:phase2_poc_metric_selection}

The previous sections compared the candidate metrics under C1--C4 and then checked smart-contract feasibility.
This section records the final selection decision for the baseline proof-of-concept. The selected metric is not
chosen because it is the easiest metric to compute. It is chosen because it has the strongest overall trade-off
between reputation meaning, stability, gaming resistance, long-term suitability, and feasibility in the intended
DESTINO-aligned hybrid architecture.

\subsection{Selected metric}
\label{sec:selected_metric_definition}

The selected PoC metric is:

\begin{quote}
\textbf{M4 developer churn ownership share}: the share of human-only churn in a frozen release window that is
attributed to one canonical developer.
\end{quote}

The concrete signal identifier is \texttt{M4\_developer\_share\_churn}. It is computed as:

\[
\texttt{share\_churn} =
\frac{\texttt{developer\_human\_churn}}
{\texttt{window\_total\_human\_churn}}
\]

For example, if a release window contains 10,000 human-churn lines and one developer is responsible for
2,500 of those churn lines, that developer's ownership share for that window is 0.25, or 25\%. The selection specifies
fixed-point integer storage rather than floating point, a rule implemented in Phase III. Let $n$ be developer churn and $d$ be
total human churn. Policy version 1.0.0 defines:

\[
\texttt{share\_churn\_ppm} =
\left\lfloor
\frac{n\times 1{,}000{,}000 + \left\lfloor d/2 \right\rfloor}{d}
\right\rfloor, \qquad d>0.
\]

This encoding rounds to the nearest ppm with exact ties upward for nonnegative integers. It avoids ambiguous
floating-point behavior and is implemented with integer arithmetic in both verifier and chaincode. Independently
rounded developer ppm values are not required to sum to exactly 1,000,000.

\subsection{Selection artifacts}
\label{sec:poc_selection_artifacts}

The selection is frozen in dedicated Phase II artifacts so the decision is reproducible and not only written in prose.

\begin{table}[H]
\centering
\scriptsize
\renewcommand{\arraystretch}{1.08}
\begin{tabularx}{\textwidth}{p{5.7cm}rX}
\toprule
\textbf{Artifact} & \textbf{Rows} & \textbf{Role} \\
\midrule
\path{out/poc_metric_selection_all.csv} & 1 & Final selected metric family, concrete signal, scores, and source tables. \\
\path{out/poc_metric_selection_rationale_all.csv} & 6 & Written reasons supporting the M4 selection. \\
\path{out/poc_metric_analysis_policy_all.csv} & 8 & Formula, population rule, encoding rule, evidence-hash rule, and supporting checks. \\
\path{out/poc_metric_record_fields_all.csv} & 13 & Selection-stage compact fields for the PoC metric record. \\
\path{out/poc_metric_supporting_signals_all.csv} & 5 & Signals retained as context or audit support around the selected metric. \\
\path{logs/poc_metric_selection_20260708.log} & 1 & Execution log for the metric-selection freeze. \\
\bottomrule
\end{tabularx}
\caption{Final PoC metric-selection artifacts. Paths are relative to the Phase II analysis directory.}
\label{tab:poc_selection_artifacts}
\end{table}

\subsection{Why M4 is selected}
\label{sec:why_m4_selected}

The selected metric comes from the M4 ownership/concentration family. The decision-input table ranked M4 first
among eligible candidates after combining data-derived C1 stability evidence, structured C2--C4 author assessments,
and smart-contract feasibility. Its
combined readiness score is 8.075. M3 activity cadence remains a strong alternative, with combined readiness
7.750, and M2 churn remains useful supporting evidence, with combined readiness 7.100.

\begin{table}[H]
\centering
\scriptsize
\renewcommand{\arraystretch}{1.08}
\begin{tabularx}{\textwidth}{p{1.0cm}p{3.5cm}cccccX}
\toprule
\textbf{ID} & \textbf{Metric} & \textbf{C1} & \textbf{C2} & \textbf{C3} & \textbf{C4} & \textbf{Feas.} & \textbf{Decision role} \\
\midrule
M4 & Ownership/concentration & 8.5 & 7 & 9 & 8 & 8.0 & Selected for the baseline PoC. \\
M3 & Activity cadence & 6.0 & 9 & 6 & 8 & 8.5 & Eligible alternative; useful as supporting activity context. \\
M2 & Churn & 4.0 & 8 & 7 & 7 & 8.0 & Eligible alternative; useful as supporting effort context. \\
M1 & Snapshot size & 10.0 & 8 & 5 & 7 & 9.0 & Not selected because it is context only. \\
M5 & Responsiveness & 0.0 & 6 & 3 & 4 & 5.0 & Not selected because current event evidence is absent. \\
\bottomrule
\end{tabularx}
\caption{Selection decision inputs for the final PoC metric.}
\label{tab:poc_selection_decision_inputs}
\end{table}

The main reasons are:
\begin{itemize}[leftmargin=*]
  \item \textbf{Reputation meaning:} M4 measures relative responsibility in a release window, not only activity volume.
  \item \textbf{Gaming-resistance assessment:} M4 received 9/10 under the structured C3 rubric because sustained ownership was judged harder to fake than simple commit splitting; this is not an empirical attack experiment.
  \item \textbf{Long-term-suitability assessment:} M4 received 8/10 under the structured C4 rubric because relative shares remain interpretable across project sizes; the current dataset does not establish universal long-term behavior.
  \item \textbf{Feasibility:} M4 scored 8.0/10 under the smart-contract feasibility rubric and triggered no exclusion condition.
  \item \textbf{Implementation boundary:} M4 can be computed off-chain and stored as a compact fixed-point record with an evidence hash.
\end{itemize}

\subsection{Analysis policy}
\label{sec:selected_metric_policy}

The selected metric must be interpreted with the same policies used in Phase I and Phase II:
\begin{itemize}[leftmargin=*]
  \item the release window must be one of the frozen \texttt{windows\_all.csv} intervals and its start commit must be an ancestor of its end commit;
  \item the developer identity must use the canonicalized author identity policy;
  \item the primary population must be human-only, with bot rows excluded from the selected metric value;
  \item the numerator is the developer's human-only churn in the window;
  \item the denominator is the total human-only churn in the same window;
  \item the fixed-point value uses the deterministic integer formula above, with scale 1,000,000 and ties rounded upward;
  \item the evidence packet must retain the analysis policy, source identifiers, numerator, denominator, encoded value, and supporting audit fields.
\end{itemize}

This policy is important because the value alone is not enough. A value of 250,000 only means 25\% ownership if
the verifier also knows the exact window, population rule, identity rule, denominator, and encoding scale.

\subsection{PoC record fields}
\label{sec:selected_metric_record_fields}

At selection time, the baseline PoC record was defined to store compact information and commitments rather than raw repository
data. Table~\ref{tab:selected_metric_record_fields} reproduces the 13-field selection artifact. Phase III refines this concept
into the exact 19-field ledger submission and 26-field immutable record schemas, including decimal-string transport for the
numerator and denominator before direct \texttt{BigInt} parsing.

\begin{table}[H]
\centering
\scriptsize
\renewcommand{\arraystretch}{1.08}
\begin{tabularx}{\textwidth}{p{3.8cm}p{2.1cm}X}
\toprule
\textbf{Field} & \textbf{Type idea} & \textbf{Role} \\
\midrule
\path{project_id} & string/bytes32 & Frozen project identifier. \\
\path{window_id} & string/bytes32 & Frozen release-window identifier. \\
\path{developer_id_hash} & bytes32 & Hash of the canonical developer identifier. \\
\path{metric_id} & string/bytes32 & Selection identifier \path{M4_developer_share_churn}. \\
\path{metric_value_ppm} & uint32/uint64 & Ownership share scaled by 1,000,000. \\
\path{numerator_churn} & uint64 & Developer human-only churn. \\
\path{denominator_churn} & uint64 & Positive total human-only window churn. \\
\path{from_commit_hash} & bytes20/bytes32 & Frozen start commit hash. \\
\path{to_commit_hash} & bytes20/bytes32 & Frozen end commit hash. \\
\path{analysis_policy_hash} & bytes32 & Hash of extraction, identity, bot, and encoding policy. \\
\path{evidence_hash} & bytes32 & Hash of the replayable off-chain evidence packet. \\
\path{storage_persistence} & string/enum & Declared persistence status or reference. \\
\path{recorded_at} & uint64 & Timestamp or block-time metadata. \\
\bottomrule
\end{tabularx}
\caption{Selection-stage compact PoC record fields for the selected M4 metric.}
\label{tab:selected_metric_record_fields}
\end{table}

\subsection{Supporting signals retained}
\label{sec:selected_metric_supporting_signals}

Selecting M4 does not mean the other implemented metrics are discarded. They remain useful as context and audit
support:
\begin{itemize}[leftmargin=*]
  \item M1 endpoint KLOC remains project-size context and normalization support.
  \item M2 developer churn shows the absolute amount of change behind an ownership share.
  \item M3 activity cadence helps distinguish sustained activity from one isolated ownership spike.
  \item M4 top-1/top-3/Gini summaries show whether the whole window is unusually concentrated.
  \item Bot summaries document whether automation affected the surrounding activity.
\end{itemize}

The selected PoC metric is therefore a primary record value, not the only evidence used to interpret the record.

\subsection{Decision boundary}
\label{sec:selected_metric_decision_boundary}

This section completes the metric-selection step. The post-hoc ancestry check in
Section~\ref{sec:posthoc_ancestry_validation} and the deterministic audit in Section~\ref{sec:phase3_poc_design}
retain M4 for the baseline PoC without claiming universal superiority. The metric definition must not change unless
the selection gate and policy version are reopened. The verifier/chaincode pair submits one replayable
\texttt{M4\_DEVELOPER\_CHURN\_SHARE} record from an eligible ancestral window. Its off-chain, local-contract, and
network results are reported in Sections~\ref{sec:phase3_offchain_verifier},~\ref{sec:phase3_chaincode_unit},
and~\ref{sec:phase3_fabric_network}.
